\section{The Magistrate on the Bishop's Chair}

In late autumn of the year 374---or, on an older chronological reckoning, 373---a provincial governor walked into the cathedral basilica of Milan to keep the peace at an episcopal election and walked out as the city's bishop-elect. He had not yet received baptism. Eight days later, baptized by a Catholic hand and hurried through the clerical orders, Ambrose sat on the episcopal throne of the Western imperial capital. Years afterward, writing to his own clergy, he still remembered the shock of the reversal: ``I was carried off from the judgment seat, and the garb of office, to enter on the priesthood, and began to teach you, what I myself had not yet learnt'' (\work{De Officiis} I.1.4). Within twenty-three years, that sudden convert from the Roman bench had helped secure the Nicene confession across the Latin West, shaped the mind of Augustine of Hippo, given Western congregations their first enduring metrical hymns, recast Cicero's civic ethics into a manual for Christian ministers, and forced Roman emperors to confront the boundary between civil power and the altar of God.

Four questions govern the life that follows. First, how did a Roman aristocrat trained in rhetoric and imperial law, rather than in a monastic school or theological academy, become one of the four original Latin Doctors of the Church? Second, what emerges when Ambrose's own contemporary letters, homilies, and treatises---together with the eyewitness testimony of Basil of Caesarea, Augustine of Hippo, and Quintus Aurelius Symmachus---are read apart from the fifth-century hagiography of Paulinus of Milan and the dramatic embroideries of Eastern historians? Third, how should a candid history weigh both the greatness and the hard edges of his public episcopate: his resistance to an Arian court over the Milanese basilicas, his debate with Symmachus over the Altar of Victory, his pastoral rebuke of Emperor Theodosius after the slaughter at Thessalonica, and his far more troubling intervention over the burned synagogue at Callinicum? Fourth, how did the historical pastor of fourth-century Milan become the Patron of the Ambrosian Church, credited by later devotion with a whole liturgical rite, scores of hymns, and the alternating composition of the \work{Te Deum}?

Milan in the late fourth century was no quiet provincial diocese. Since the tetrarchy, the northern Italian city had served repeatedly as an imperial residence---closer than Rome to the Danube and Rhine frontiers, and home to the court, the praetorian prefecture of Italy, the military high command, and a restless garrison that included Gothic federates. When Ambrose assumed the see, the Nicene settlement of 325 remained fiercely contested. For nearly two decades, ever since the Emperor Constantius~II exiled the Catholic bishop Dionysius at the synod of Milan in 355, the Milanese cathedral had been held by Auxentius of Cappadocia, a Homoian bishop who rejected the Nicene confession that the Son is of one substance with the Father. Around the court stood pagan Roman senators who still revered the ancestral rites, Homoian courtiers who viewed the imperial palace as the arbiter of church property, and Catholic congregations hungry for orthodox teaching and steady pastoral care. Into that divided capital stepped a Roman magistrate whose life would test, again and again, what belongs to Caesar and what belongs to God.
